\makereporttitle{2025}{martie}

\makereportpreamble

\vspace{0.5em}


În martie am continuat să explorez apariția haosului în modelul cuadrupol bosonic pentru oscilații nucleare. Am identificat noi comportamente interesante în sistemul studiat.

În contextul sistemului dinamic obținut în limita semiclasică, am reușit să calculez valoarea maximă a exponentului Lyapunov pentru mai multe condiții inițiale și diferite valori ale parametrilor \(B\), \(D\) și a energiei totale. Am potrivit o funcție exponențială pentru valoarea exponentului Lyapunov în funcție de valoarea parametrului \(B\) în regiunea de după valoarea critică \(B_0\).

\begin{figure}[htbp]
    \centering
    \includegraphics[width=0.8\linewidth]{reports/2025/images/3/lyapunov_exponent_power_law}
\end{figure}

În ceea ce privește nivelele de energie ale sistemului cuantic, am reușit în sfârșit să înțeleg cum pot clasifica valorile/vectorii proprii ai Hamiltonianului scris într-o bază 2D de oscilatori armonici cuantici, în funcție de reprezentările ireductibile ale grupului \(C_{3v}\) de care aparțin. Încă nu am reușit să reproduc rezultatele din lucrarea de cercetare ``\textit{Fingerprints of classical chaotic dynamics in quantum behavior}'', de S.~Micluță-Câmpeanu, M.C.~Raportaru, A.I.~Nicolin și V.~Băran.

Pe partea educațională, am participat la prelegerile despre ``\emph{Standard model of fundamental interactions}'', ținute de prof.~dr.~Emilian Dudaș, venit în vizită la Facultatea de Fizică a UB. Acestea m-au ajutat să înțeleg modelul standard al particulelor elementare și ce fel de probleme se cercetează în fizica teoretică modernă. De asemenea, am continuat să particip la diferite seminare științifice organizate în cadrul Institutului de Matematică ``Simion Stoilow'' al Academiei Române (inclusiv despre cosmologie și geometrie riemanniană) și am parcurs mai multe lecții din cursul online de ``\textit{Mathematical and computational methods for Physics}'', de la Georgetown University.

De asemenea, luna aceasta am progresat în implementarea proiectului DANUBIUS, configurând și lansând o versiune inițială a viitorului \href{https://data.danubius-ri.eu/}{portal de date} pentru consorțiu. Aceasta este o instanță a aplicației \textit{open source} \href{https://ckan.org/}{CKAN}, instalată pe infrastructura din ACC. Am participat la conferința \href{https://uefiscdi.gov.ro/bridging-disciplines-advancing-research-data-management-for-collaborative-innovation-conference}{\textit{Bridging Disciplines: Advancing Research Data Management for Collaborative Innovation}}, organizată pe 20--21 martie de \href{https://uefiscdi.gov.ro/}{UEFISCDI}. În cadrul acesteia, am aflat despre cele mai noi direcții la nivel european în ceea ce privește \emph{(research) data stewardship} și am avut oportunitatea să cunosc persoane cu roluri variate și o experiență deosebită în domeniu.


\textbf{Observațiile conducătorului științific}:

\vspace{3em}

\makereportpostamble{27 martie 2025}
